% Full-model signed product identities, derived from the original generator.
% No irrationality or transcendence conclusion is asserted.

\section{A signed hypergeometric product martingale}
\label{sec:bipartite-product-martingale}

Consider the original free friend tree, started from one edge, or the
one-stub inert tree, started from a real root of full degree two with one
real leaf and one inert stub.  Every real vertex has a rate-one source
clock and selects a uniformly random full neighbor.  A selection of the
inert stub has no effect.  Choose a consistent bipartition
$\sigma_v\in\{1,-1\}$ of the real tree; in the inert case take the root
to have sign $1$.  A new leaf attached to $w$ has sign $-\sigma_w$.
Degrees $d_v$ always include the inert stub, if present, but every sum
or product over vertices below includes only real vertices.

For a state $T$, put
\[
 h_T(w)=\sum_{u\sim w\text{ real}}\frac1{d_u},\qquad
 (\mathcal L f)(T)=\sum_{w\text{ real}}h_T(w)
       \bigl(f(T+w)-f(T)\bigr),
\]
where $T+w$ denotes the attachment of one new leaf at $w$.  Thus the
proper birth rate is $n$ in the free model and $n-1/d_0$ in the inert
model, where $n$ is the real population and $0$ is the inert root.

\begin{proposition}[Polynomial harmonic family]
\label{prop:bipartite-product}
For every complex parameter $z$, the finite-state polynomial
\begin{equation}
 P_z(T):=\prod_{v\text{ real}}\prod_{j=1}^{d_v-1}
       \left(1+\frac{\sigma_vz}{j}\right)
   =\prod_{v\text{ real}}
       \frac{(1+\sigma_vz)_{d_v-1}}{(d_v-1)!}
 \label{eq:bipartite-P}
\end{equation}
satisfies $\mathcal LP_z=0$ as a polynomial identity in $z$.
If $T_n$ is the embedded process at real population $n$, then
\begin{equation}
 \mathbb E P_z(T_n)=
 \begin{cases}
  1,&\text{free initial edge},\\
  1+z,&\text{inert newborn state}.
 \end{cases}
 \label{eq:bipartite-normalization}
\end{equation}
Every fixed coefficient $[z^k]P_z(T_t)$ is a true martingale in continuous
time.  The whole product $P_z(T_t)$ is a true martingale for real
$-1\le z\le1$.
\end{proposition}

\begin{proof}
The newborn has degree one, so its factor in \eqref{eq:bipartite-P} is
one.  Only the target's factor changes.  Consequently
\begin{equation}
 P_z(T+w)=P_z(T)\left(1+\frac{\sigma_wz}{d_w}\right).
 \label{eq:bipartite-birth-ratio}
\end{equation}
This is a polynomial identity; it does not divide by $P_z(T)$ and remains
valid at its zeros.  We therefore have
\begin{align*}
 (\mathcal LP_z)(T)
 &=zP_z(T)\sum_w\frac{\sigma_w}{d_w}
                 \sum_{u\sim w\text{ real}}\frac1{d_u}\\
 &=zP_z(T)\sum_{\{u,v\}\text{ real edge}}
                  \frac{\sigma_u+\sigma_v}{d_ud_v}=0.
\end{align*}
No source is assigned to the inert stub, so there is no unmatched
boundary term.  Dividing the generator identity by the positive proper
birth rate gives the embedded martingale identity.  There are finitely
many states at each fixed population.  Its expectation is therefore
well defined for every complex $z$, and the initial products are
respectively $1$ and $1+z$.

For the continuous-time coefficient assertion, the number of linear
factors in $P_z$ is $n-2$ in the free case and $n-1$ in the inert case.
Every factor's nonconstant coefficient has absolute value at most one.
It follows that
\[
 \bigl|[z^k]P_z(T)\bigr|\le {n-1\choose k}\le n^k.
\]
Both real populations are dominated by the rate-one Yule process
started from two vertices.  This process has finite moments of every
fixed order at each finite time.  Stop at population $M$, apply the
finite-state martingale identity, and let $M$ tend to infinity.
The displayed bound, monotonicity of the population, and dominated
convergence give the true martingale identity, including conditional
expectations, for each fixed coefficient.

For real $|z|<1$, $P_z$ is strictly positive.  Its Doob transform has
attachment rates
\begin{equation}
 h_T^{(z)}(w)
    =h_T(w)\left(1+\frac{\sigma_wz}{d_w}\right).
 \label{eq:bipartite-doob-rates}
\end{equation}
They are nonnegative, and their sum equals the original proper birth
rate by the same edge cancellation.  This transformed pure-birth
process is nonexplosive, since its total rate at population $n$ is at
most $n$.  On every finite path its likelihood ratio to the original
process is $P_z(T_t)/P_z(T_0)$: the waiting rates agree and the jump-rate
ratios telescope.  Summing over all finite paths gives expectation one
for this ratio.  The Markov property proves the martingale assertion.

At $z=\pm1$, restrict the same argument to states where $P_z>0$ when
the initial value is positive.  Transitions to states with $P_z=0$ have
transformed rate zero; all other rates are nonnegative and their sum
is unchanged.  Once the original process has $P_z=0$, the polynomial
birth identity shows that it remains zero.  For the inert initial
state at $z=-1$, the product is identically zero throughout, since the
positive root always has degree at least two.  This treats all endpoints.
No assertion of integrability of the whole product is needed here for
other complex parameters.
\end{proof}

The family in \eqref{eq:bipartite-P} is hypergeometric in every individual
degree: the ratio of consecutive degree factors is
$(d+\sigma_vz)/d$.  This is an identity for the full original process,
not a replacement by independent degree processes.

\begin{proposition}[A signed martingale with the core-growth exponent]
\label{prop:bipartite-variance}
Write
\[
 H_m=\sum_{j=1}^m\frac1j,\qquad
 H_m^{(2)}=\sum_{j=1}^m\frac1{j^2},\qquad H_0=H_0^{(2)}=0,
\]
and define
\begin{equation}
 J(T)=\sum_v\sigma_vH_{d_v-1},\qquad
 S_2(T)=\sum_vH_{d_v-1}^{(2)}.
 \label{eq:bipartite-J-S2}
\end{equation}
Then $J(T_t)$ and $J(T_t)^2-S_2(T_t)$ are true martingales.
For the free edge initial state, $\mathbb EJ(T_t)=0$; for the inert
newborn initial state, $\mathbb EJ(T_t)=1$.  In both models,
\begin{equation}
 \mathbb EJ(T_t)^2=\mathbb ES_2(T_t),\qquad
 \mathbb EC(T_t)\le\mathbb EJ(T_t)^2
       \le\zeta(2)\mathbb EC(T_t).
 \label{eq:bipartite-variance-comparison}
\end{equation}
The same identities hold at every fixed real population.
In particular, using the known growth exponent of the original free
tree,
\begin{equation}
 \boxed{\mu=\lim_{t\to\infty}\frac{\log\mathbb EJ(T_t)^2}{t}
       =\lim_{n\to\infty}
          \frac{\log\mathbb EJ(T_n)^2}{\log n}.}
 \label{eq:bipartite-mu-variance}
\end{equation}
The second equality invokes the known expected-population exponent
$\mathbb EC(T_n)=n^{\mu+o(1)}$; the first invokes its continuous-time
counterpart.  The two martingale identities themselves require neither
growth theorem.
\end{proposition}

\begin{proof}
The first two nonconstant coefficients of the product are
\[
 P_z(T)=1+zJ(T)+\frac{z^2}{2}\bigl(J(T)^2-S_2(T)\bigr)+O(z^3).
\]
Apply Proposition~\ref{prop:bipartite-product}.  At the free initial
edge, $(J,S_2)=(0,0)$, while at the inert initial state
$(J,S_2)=(1,1)$.  Thus $J^2-S_2$ has expectation zero in either model.
Every leaf contributes zero to $S_2$, and every core vertex contributes
a number between $1$ and $\zeta(2)$.  This proves the comparison, which
then transfers either already known growth exponent to $\mathbb EJ^2$.
\end{proof}

There is also a pathwise quadratic-variation identity.  At a birth
attached to $w$,
\[
 \Delta J=\frac{\sigma_w}{d_w},\qquad
 \Delta S_2=\frac1{d_w^2}=(\Delta J)^2.
\]
Consequently the optional quadratic variation satisfies
\begin{equation}
 [J]_t=S_2(T_t)-S_2(T_0),
 \qquad
 \langle J\rangle_t=\int_0^t\sum_w\frac{h_{T_u}(w)}{d_w(u)^2}\,du.
 \label{eq:bipartite-brackets}
\end{equation}
In particular $|\Delta J|\le1$.  The optional and predictable brackets
are different objects; only their expectations need agree.

\begin{proposition}[Signed integer specialization]
\label{prop:bipartite-integer}
Let $k\ge1$ be an integer and let $N_+$ denote the number of positive
real vertices.  For either the free tree or the root-positive inert
tree,
\begin{equation}
 P_k(T)=(-1)^{N_+-1}
   \prod_{\sigma_v=1}{k+d_v-1\choose d_v-1}
   \prod_{\sigma_v=-1}{k-1\choose d_v-1},
 \label{eq:bipartite-binomial}
\end{equation}
where the second binomial coefficient is zero if $d_v>k$.
Thus $P_k$ is a signed integer observable, supported on trees in which
every negative vertex has degree at most $k$.  Its expectation at
fixed population is exactly $1$ in the free tree and $k+1$ in the inert
tree.

For the free tree and $n\ge3$,
\begin{equation}
 \mathbb P(T_n\text{ is a star with positive center})=\frac1{n-1},
 \qquad \mathbb P(T_n\text{ is a star})=\frac2{n-1}.
 \label{eq:bipartite-free-star}
\end{equation}
For the inert tree and $n\ge2$,
\begin{equation}
 \mathbb P(T_n\text{ is a star centered at its real root})=\frac2n.
 \label{eq:bipartite-inert-star}
\end{equation}
Here the inert root's full degree in such a star is $n$.
\end{proposition}

\begin{proof}
The positive degree factor is
$(1+k)_{d-1}/(d-1)!={k+d-1\choose d-1}$, while the negative factor is
$(1-k)_{d-1}/(d-1)!=(-1)^{d-1}{k-1\choose d-1}$.
Every real edge has exactly one negative endpoint.  The inert stub, if
present, is attached to a positive vertex.  Hence in both models
\[
 \sum_{\sigma_v=-1}(d_v-1)=(n-1)-N_-=N_+-1,
\]
which proves the sign in \eqref{eq:bipartite-binomial}.  The expectation
follows from \eqref{eq:bipartite-normalization}.

At $k=1$, all negative vertices must be leaves on the support of $P_1$.
A connected tree with this property has exactly one positive vertex,
unless it is the initial free edge.  The supporting state is therefore
a positive-centered star.  Its product is $n-1$ in the free model.
Color symmetry gives the second free identity.  In the inert model,
the positive root is already a core vertex, so the only possible
positive center is this root.  Its product is $n$, and its initial
expectation is $2$, proving \eqref{eq:bipartite-inert-star}.
\end{proof}

% Actual free friend-tree identities; no arithmetic classification is asserted.
\subsection{Core and color insertions into the signed product}
\label{bpi:section}

Work with the original free tree, its fixed bipartition
$\sigma_v\in\{+1,-1\}$, and attachment rates
$h_T(w)=\sum_{u\sim w}d_u^{-1}$.  Write
\[
 P_z(T)=\prod_v\prod_{j=1}^{d_v-1}
        \left(1+\frac{\sigma_vz}{j}\right),
 \qquad Q_{z,x}(T)=x^{C(T)}P_z(T),
\]
and define the leaf-promotion and edge observables
\begin{align*}
 \Lambda&=\sum_{w:\,d_w=1}h_T(w),&
 \Lambda_\sigma&=\sum_{w:\,d_w=1}\sigma_wh_T(w),\\
 W&=\sum_{\{u,v\}\in E(T)}\frac1{d_ud_v},&
 D&=\sum_v\sigma_v=N_+-N_-.
\end{align*}

\begin{proposition}[Exact insertion identities]
\label{bpi:identities}
For every finite free tree with at least two vertices, the following are
polynomial identities in $x,z\in\mathbb C$:
\begin{align}
 \mathcal LQ_{z,x}
   &=(x-1)Q_{z,x}(\Lambda+z\Lambda_\sigma),
       \label{bpi:fugacity}\\
 \mathcal L(P_zC)
   &=P_z(\Lambda+z\Lambda_\sigma),
       \label{bpi:core}\\
 \mathcal L(P_zD)
   &=P_z(D-2zW).
       \label{bpi:color}
\end{align}
For the edge-start embedded chain, set
\[
 a_n(z)=\mathbb E[P_z(T_n)C(T_n)],\qquad
 d_n(z)=\mathbb E[P_z(T_n)D(T_n)].
\]
Then, for every $n\geq2$ and every complex $z$,
\begin{align}
 n\bigl(a_{n+1}(z)-a_n(z)\bigr)
   &=\mathbb E\!\left[P_z(T_n)
                (\Lambda(T_n)+z\Lambda_\sigma(T_n))\right],
       \label{bpi:population-core}\\
 n\bigl(d_{n+1}(z)-d_n(z)\bigr)
   &=d_n(z)-2z\mathbb E[P_z(T_n)W(T_n)].
       \label{bpi:population-color}
\end{align}
Here $a_2(z)=d_2(z)=0$ and $a_n(0)=\mathbb EC(T_n)$.
\end{proposition}

\begin{proof}
For attachment at $w$, put $\delta_w=\mathbf1_{\{d_w=1\}}$.
The new leaf has degree one and sign $-\sigma_w$.  Thus, as polynomial
identities without division by $P_z$ or by $x$,
\[
 P_z(T+w)=P_z(T)\left(1+\frac{\sigma_wz}{d_w}\right),
 \qquad C(T+w)=C(T)+\delta_w.
\]
Consequently
\[
 Q_{z,x}(T+w)-Q_{z,x}(T)
 =Q_{z,x}(T)\left[
   \frac{\sigma_wz}{d_w}
   +(x-1)\delta_w(1+\sigma_wz)\right].
\]
The first term vanishes after summation against the rates, since
\[
 \sum_w\frac{\sigma_wh_T(w)}{d_w}
 =\sum_{\{u,v\}\in E(T)}
       \frac{\sigma_u+\sigma_v}{d_ud_v}=0.
\]
This proves \eqref{bpi:fugacity}.  Differentiation in $x$ at $x=1$
proves \eqref{bpi:core}.

A birth at $w$ changes $D$ by $-\sigma_w$, and hence
\[
 \mathcal LD=-\sum_w\sigma_wh_T(w)=\sum_u\sigma_u=D.
\]
The product rule for the jump generator, together with
$\mathcal LP_z=0$, now gives
\[
 \mathcal L(P_zD)
 =P_zD-zP_z\sum_w\frac{h_T(w)}{d_w}
 =P_z(D-2zW),
\]
because each undirected edge occurs twice in the last sum.
Finally, the total birth rate is $n$, so conditioning on $T_n$ gives
\eqref{bpi:population-core} and \eqref{bpi:population-color}.
Every population level has finitely many tree states, making these
expectations finite polynomials for all complex parameter values.
\end{proof}

These identities remove the nonleaf terms but retain the actual neighbor
dependence in $\Lambda\pm\Lambda_\sigma$.  They do not close a finite
scalar recurrence or establish irrationality.  Their all-complex domain
is a fixed-state or fixed-population statement; no all-complex
continuous-time integrability assertion is being made.


\paragraph{Arithmetic scope and unresolved step.}
The cancellation in Proposition~\ref{prop:bipartite-product} is exact,
signed, and statewise.  The integer specializations in
\eqref{eq:bipartite-binomial} require no clearing of degree denominators.
More generally, if $D_n=\operatorname{lcm}(1,\ldots,n-1)$, then
$D_n^k[z^k]P_z(T)$ is an integer at every state of real population $n$.
This follows directly by expanding the product into elementary
symmetric functions of numbers $\sigma_v/j$, whose denominators all
divide $D_n$.  In particular $D_nJ(T)$ and $D_n^2S_2(T)$ are integers.
These are state-observable denominator bounds, not bounds for their
expectations: the state probabilities introduce additional denominators.
The cancellation occurs before any expectation or denominator clearing.  However,
$z$ is a new deformation parameter, not the critical exponent $\mu$.
The identities $\mathbb EP_k=1$ or $k+1$ are probability-weighted
normalizations, and do not make these expectations sums of integers.
They do not currently supply a nonzero integer expression tending to
zero when $\mu$ is rational.  Likewise, the second-moment realization
\eqref{eq:bipartite-mu-variance} identifies the same growth exponent but
does not constrain its arithmetic type.  A new quantitative bridge
from the signed product identities to the critical growth equation is
still required for irrationality.


% This comparison process is not the original friend tree.
\subsection{What the signed product martingales and color process do not determine}

The following comparison process preserves both the full family of signed
product martingales and the exact two-color population process, but has
rational core-growth exponent.  In particular these two structures, even
together with arbitrarily many original finite coefficients, cannot by
themselves classify the arithmetic nature of the original exponent.

\begin{proposition}[A degree-weighted comparison process]\label{bpm:comparison}
Fix an integer $K\geq3$.  Start from an edge with its fixed bipartition
$\sigma_v\in\{+1,-1\}$, use the original friend rule through population
$K$, and thereafter attach a new leaf at $v$ at rate
\[
 \lambda_v(T)=\frac{N_{-\sigma_v}(T)d_v(T)}{n-1},
 \qquad n=|T|\geq K.
\]
The new vertex receives sign $-\sigma_v$.  This process has rational
transition probabilities and total birth rate $n$.  For every complex
$z$, the polynomial
\[
 P_z(T)=\prod_v\prod_{j=1}^{d_v-1}\left(1+\frac{\sigma_vz}{j}\right)
\]
has $\mathbb E P_z(T_n)=1$ at every population.  Its two color populations
in continuous time are independent rate-one Yule processes, each starting
from one vertex, exactly as for the original friend process.  Nevertheless,
if $a_n^*=\mathbb E C(T_n)$ for this comparison process, then
\[
 \frac n6-O_K(1)\leq a_n^*\leq n-2,
 \qquad a_n^*=n^{1+o(1)}.
\]
All population coefficients through $K$ and all edge-start generator
derivatives through order $K-2$ agree with those of the original process.
\end{proposition}

\begin{proof}
The sum of the degrees in either bipartition class is $n-1$.  Therefore
\[
 \sum_{v:\sigma_v=+1}\lambda_v=N_-,\qquad
 \sum_{v:\sigma_v=-1}\lambda_v=N_+,
\]
and the total rate is $n$.  In particular the continuous-time count pair
has generator
\[
 (\mathcal Af)(i,j)
 =i\bigl(f(i+1,j)-f(i,j)\bigr)
  +j\bigl(f(i,j+1)-f(i,j)\bigr),
\]
which is the generator of two independent Yule processes. For the
original rule, each newborn has its source vertex's color, so
the same projection holds before the switch and hence throughout.
This is the continuous-time embedding of the known P\'olya color urn
\cite[Theorem 3.1]{CanningsJordan}.
Equivalently, at population $n$, $N_+$ is uniform on
$\{1,\ldots,n-1\}$.

Attaching at $v$ changes $P_z$ by the polynomial identity
\[
 P_z(T+v)-P_z(T)=\frac{\sigma_vz}{d_v}P_z(T).
\]
This identity remains valid at roots of $P_z$; no division by $P_z$ is
being used.  For the modified rates,
\[
 \sum_v\frac{\sigma_v\lambda_v}{d_v}
 =\frac{N_+N_--N_-N_+}{n-1}=0.
\]
Consequently $\mathcal L^*P_z=0$.  The original generator has the same
identity.  Conditioning on the finite next-birth distribution therefore
gives $\mathbb E P_z(T_n)=P_z(T_2)=1$; no unbounded-time optional-stopping
argument is needed.

Let $L_+,L_-$ be the numbers of leaves of each sign, $C$ the number of
nonleaves, and $m=\min(N_+,N_-)$.  For $n\geq K$, the exact promotion
rate and its lower bound are
\[
 \mathcal L^*C
 =\frac{N_-L_++N_+L_-}{n-1}
 \geq\frac{m(n-C)}{n-1}.
\]
Since $m\leq n/2$ deterministically and its uniform-color expectation
satisfies $\mathbb E m\geq n/4$, division by the total birth rate $n$
gives
\[
 a_{n+1}^*-a_n^*
 \geq\frac{n}{4(n-1)}-\frac{a_n^*}{2(n-1)}.
\]
Writing $b_n=a_n^*-n/6$, this becomes
\[
 b_{n+1}\geq
 \left(1-\frac1{2(n-1)}\right)b_n+\frac1{6(n-1)}.
\]
It follows inductively that $b_n\geq\min(b_K,0)$ for every $n\geq K$.
The upper bound $C\leq n-2$ is deterministic.  These bounds prove the
claimed exponent.  Agreement of the finite coefficients follows from
using exactly the original transitions below population $K$.  A generator
word of order at most $K-2$, started at the edge and applied to $C$, uses
no modified transition, proving agreement of the stated derivatives.
\end{proof}

At an integer $z=k\geq1$ these same comparison processes consequently
satisfy every signed binomial identity
\[
 \mathbb E\left[
 (-1)^{N_+-1}
 \prod_{\sigma_v=+1}\binom{k+d_v-1}{d_v-1}
 \prod_{\sigma_v=-1}\binom{k-1}{d_v-1}
 \right]=1,
\]
with the usual convention that the second binomial is zero when
$d_v>k$.  The comparison does not identify the original exponent, alter
the original process, or refute an arithmetic theorem using additional
information specific to the friend rates
$h_T(v)=\sum_{u\sim v}d_u^{-1}$.  It isolates the missing ingredient:
the local neighbor dependence must enter any such theorem, beyond the
color balance and the product-martingale cancellation.
